\documentclass[11pt,a4paper]{article}

\usepackage[utf8]{inputenc}
\usepackage[T1]{fontenc}
\usepackage[spanish]{babel}
\usepackage{geometry}
\usepackage{longtable}
\usepackage{booktabs}
\usepackage{array}
\usepackage{multirow}
\usepackage{enumitem}
\usepackage{tabularx}
\usepackage{graphicx}
\usepackage{tikz}
\usepackage{forest}
\usepackage{hyperref}
\usepackage{titlesec}
\usepackage{xcolor}

\geometry{margin=2.5cm}

\title{Síntesis y Reconstrucción Taxonómica del Plan Integrado del Área de Matemáticas}
\author{}
\date{}

\begin{document}
	
	\maketitle
	
	\section*{Introducción}
	
	El presente documento consolida las taxonomías, clasificaciones, sistemas de categorización, estructuras jerárquicas, dominios conceptuales, etapas de desarrollo, componentes curriculares y marcos organizativos identificados en el Plan Integrado del Área de Matemáticas. Se reconstruyen las relaciones entre dichas estructuras manteniendo la terminología original y explicitando los vínculos conceptuales observados.
	
	\tableofcontents
	\newpage
	
	\section{Visión General de las Taxonomías Identificadas}
	
	Se identificaron los siguientes sistemas taxonómicos principales:
	
	\begin{enumerate}
		\item Sistemas Matemáticos.
		\item Tipos de Pensamiento Matemático.
		\item Subprocesos del Área de Matemáticas.
		\item Etapas de Desarrollo de la Modelación Matemática.
		\item Etapas de Desarrollo de la Comunicación Matemática.
		\item Etapas de Desarrollo del Razonamiento Matemático.
		\item Sistemas de Representación Semiótica.
		\item Actividades Cognitivas de la Semiosis.
		\item Tipos de Razonamiento.
		\item Componentes de la Comunicación Matemática.
		\item Enfoques de Caracterización de los Modelos.
		\item Clasificación de la Resolución de Problemas.
		\item Tipos de Procedimientos para Resolver Problemas.
		\item Objetivos de la Resolución de Problemas.
		\item Estructuras de Valores y Estatus en el Razonamiento.
		\item Campos Disciplinares que Nutren la Educación Matemática.
	\end{enumerate}
	
	\section{Taxonomía 1: Sistemas Matemáticos}
	
	\subsection{Descripción}
	
	Conjunto de sistemas mediante los cuales se desarrolla el pensamiento matemático.
	
	\subsection{Propósito}
	
	Organizar los objetos matemáticos y sus relaciones para favorecer el desarrollo del pensamiento matemático.
	
	\subsection{Estructura jerárquica}
	
	\begin{forest}
		for tree={
			draw,
			rounded corners,
			align=center
		}
		[Sistemas Matemáticos
		[Sistema Numérico]
		[Sistema Métrico]
		[Sistema Geométrico]
		[Sistema Algebraico]
		[Sistema Analítico y de Datos]
		]
	\end{forest}
	
	\subsection{Tabla resumen}
	
	\begin{longtable}{p{4cm}p{9cm}}
		\toprule
		Sistema & Características principales\
		\midrule
		Numérico & Números, operaciones, propiedades, relaciones y algoritmos.\
		Geométrico & Espacio, figuras, representación espacial y argumentación.\
		Métrico & Magnitudes, unidades, medición, precisión y error.\
		Algebraico & Variable, relaciones, cambios y modelación simbólica.\
		Analítico y de Datos & Datos, incertidumbre, inferencia, probabilidad y toma de decisiones.\
		\bottomrule
	\end{longtable}
	
	\subsection{Relación con otras taxonomías}
	
	Estos sistemas sustentan los tipos de pensamiento matemático y atraviesan todos los subprocesos.
	
	\section{Taxonomía 2: Tipos de Pensamiento Matemático}
	
	\subsection{Descripción}
	
	Clasificación de formas de pensamiento asociadas al desarrollo histórico de las matemáticas.
	
	\subsection{Propósito}
	
	Explicar las distintas maneras de construir conocimiento matemático.
	
	\subsection{Estructura jerárquica}
	
	\begin{forest}
		for tree={draw,rounded corners}
		[Pensamiento Matemático
		[Pensamiento Numérico]
		[Pensamiento Espacial]
		[Pensamiento Métrico]
		[Pensamiento Aleatorio]
		[Pensamiento Analítico o Variacional]
		]
	\end{forest}
	
	\subsection{Tabla resumen}
	
	\begin{longtable}{p{4cm}p{9cm}}
		\toprule
		Tipo & Asociado principalmente a\
		\midrule
		Numérico & Sistema numérico\
		Espacial & Sistema geométrico\
		Métrico & Sistema de medida\
		Aleatorio & Probabilidad e incertidumbre\
		Analítico o Variacional & Cambio y variación\
		\bottomrule
	\end{longtable}
	
	\subsection{Relación con otras taxonomías}
	
	Surgen del desarrollo de los sistemas matemáticos.
	
	\section{Taxonomía 3: Subprocesos del Área de Matemáticas}
	
	\subsection{Descripción}
	
	Procesos transversales que evidencian el desarrollo del pensamiento matemático.
	
	\subsection{Propósito}
	
	Orientar la formación matemática integral.
	
	\subsection{Estructura jerárquica}
	
	\begin{forest}
		for tree={draw,rounded corners}
		[Subprocesos
		[Modelación Matemática]
		[Comunicación Matemática]
		[Razonamiento Matemático]
		[Resolución y Formulación de Problemas]
		]
	\end{forest}
	
	\subsection{Tabla resumen}
	
	\begin{longtable}{p{5cm}p{8cm}}
		\toprule
		Subproceso & Finalidad principal\
		\midrule
		Modelación Matemática & Representar y comprender fenómenos.\
		Comunicación Matemática & Construir y compartir significados.\
		Razonamiento Matemático & Justificar, demostrar y validar.\
		Resolución y Formulación de Problemas & Resolver, extender y crear problemas.\
		\bottomrule
	\end{longtable}
	
	\subsection{Relación con otras taxonomías}
	
	Todos los subprocesos utilizan sistemas matemáticos y representaciones semióticas.
	
	\section{Taxonomía 4: Etapas de la Modelación Matemática}
	
	\subsection{Descripción}
	
	Secuencia de progresión del proceso de modelación.
	
	\subsection{Propósito}
	
	Orientar el desarrollo gradual de la modelación según el grado escolar.
	
	\subsection{Estructura jerárquica}
	
	\begin{forest}
		for tree={draw,rounded corners}
		[Modelación Matemática
		[Reconocimiento de Propiedades\1°--3°]
		[Identificación de Regularidades\4°--6°]
		[Transferencia de Problemas y Modelos\7°--9°]
		[Verificación de Modelos\10°--11°]
		]
	\end{forest}
	
	\subsection{Tabla resumen}
	
	\begin{longtable}{p{4cm}p{2cm}p{7cm}}
		\toprule
		Etapa & Grados & Énfasis\
		\midrule
		Reconocimiento de Propiedades & 1°--3° & Procedimientos aritméticos, métricos y geométricos.\
		Identificación de Regularidades & 4°--6° & Detección de patrones y procedimientos repetitivos.\
		Transferencia de Problemas y Modelos & 7°--9° & Traducción entre realidad y modelo matemático.\
		Verificación de Modelos & 10°--11° & Validación e interpretación de resultados.\
		\bottomrule
	\end{longtable}
	
	\subsection{Relación con otras taxonomías}
	
	Se articula directamente con la matematización, las representaciones semióticas y la resolución de problemas.
	
	\section{Taxonomía 5: Etapas de la Comunicación Matemática}
	
	\subsection{Descripción}
	
	Evolución de la comunicación matemática a lo largo de la escolaridad.
	
	\subsection{Propósito}
	
	Favorecer la transición desde el lenguaje cotidiano al lenguaje formal.
	
	\subsection{Estructura jerárquica}
	
	\begin{forest}
		for tree={draw,rounded corners}
		[Comunicación Matemática
		[Expresión y verbalización\1°--3°]
		[Interpretación\4°--6°]
		[Argumentación\7°--9°]
		[Demostración y formalización\10°--11°]
		]
	\end{forest}
	
	\subsection{Relación con otras taxonomías}
	
	Se apoya en las representaciones semióticas y en el razonamiento matemático.
	
	\section{Taxonomía 6: Componentes de la Comunicación Matemática}
	
	\subsection{Descripción}
	
	Elementos fundamentales del acto comunicativo en educación matemática.
	
	\subsection{Estructura jerárquica}
	
	\begin{forest}
		for tree={draw,rounded corners}
		[Comunicación Matemática
		[Saber Específico
		[Numérico--Variacional]
		[Geométrico--Métrico]
		[Aleatorio]
		]
		[Actores Educativos
		[Maestro--Maestro]
		[Maestro--Estudiante]
		[Estudiante--Estudiante]
		]
		[Escenarios Educativos]
		[Ambiente Educativo]
		]
	\end{forest}
	
	\section{Taxonomía 7: Sistemas de Representación Semiótica}
	
	\subsection{Descripción}
	
	Representaciones utilizadas para acceder a los objetos matemáticos.
	
	\subsection{Estructura jerárquica}
	
	\begin{forest}
		for tree={draw,rounded corners}
		[Representaciones Semióticas
		[Lenguaje Natural]
		[Representaciones Simbólicas]
		[Representaciones Gráficas]
		[Representaciones Pictóricas]
		[Tablas]
		[Gráficas Cartesianas]
		[Fórmulas]
		[Expresiones Analíticas]
		]
	\end{forest}
	
	\subsection{Nota}
	
	El documento menciona múltiples registros de representación; algunos aparecen explícitamente en contextos específicos y otros como ejemplos de registros.
	
	\section{Taxonomía 8: Actividades Cognitivas de la Semiosis}
	
	\subsection{Descripción}
	
	Actividades cognitivas propuestas por Raymond Duval.
	
	\subsection{Estructura jerárquica}
	
	\begin{forest}
		for tree={draw,rounded corners}
		[Semiosis
		[Formación]
		[Tratamiento]
		[Conversión]
		]
	\end{forest}
	
	\subsection{Tabla resumen}
	
	\begin{longtable}{p{4cm}p{9cm}}
		\toprule
		Actividad & Función\
		\midrule
		Formación & Construcción de representaciones.\
		Tratamiento & Transformación dentro del mismo registro.\
		Conversión & Transformación entre registros distintos.\
		\bottomrule
	\end{longtable}
	
	\section{Taxonomía 9: Criterios de Congruencia Representacional}
	
	\subsection{Descripción}
	
	Criterios utilizados para determinar la congruencia entre representaciones.
	
	\begin{enumerate}
		\item Correspondencia semántica.
		\item Univocidad semántica terminal.
		\item Conservación del orden de organización.
	\end{enumerate}
	
	\section{Taxonomía 10: Enfoques para la Caracterización de un Modelo}
	
	\subsection{Descripción}
	
	Perspectivas desde las cuales puede analizarse un modelo.
	
	\begin{forest}
		for tree={draw,rounded corners}
		[Modelo
		[Enfoque de Sistema]
		[Enfoque Causal]
		[Enfoque Dialéctico]
		[Enfoque Genético]
		]
	\end{forest}
	
	\section{Taxonomía 11: Tipos de Razonamiento}
	
	\subsection{Descripción}
	
	Clasificación explícita de formas de razonamiento.
	
	\begin{forest}
		for tree={draw,rounded corners}
		[Razonamiento
		[Deductivo]
		[Reducción al Absurdo]
		[Inductivo]
		]
	\end{forest}
	
	\subsection{Relación con otras taxonomías}
	
	Constituyen la base del subproceso de razonamiento matemático.
	
	\section{Taxonomía 12: Estructura del Razonamiento Demostrativo}
	
	\subsection{Descripción}
	
	Componentes identificados en el razonamiento demostrativo.
	
	\begin{forest}
		for tree={draw,rounded corners}
		[Razonamiento Demostrativo
		[Valores
		[Valor Epistémico]
		[Valor Lógico de Verdad]
		]
		[Estatus
		[Estatus Teórico]
		[Estatus Operatorio]
		]
		]
	\end{forest}
	
	\subsection{Fases}
	
	\begin{enumerate}
		\item Antes (Conjetura).
		\item Durante (Deducción).
		\item Después (Necesidad de la conclusión).
	\end{enumerate}
	
	\section{Taxonomía 13: Clasificación de la Resolución de Problemas}
	
	\subsection{Descripción}
	
	Clasificación basada en Stanic y Kilpatrick.
	
	\begin{forest}
		for tree={draw,rounded corners}
		[Resolución de Problemas
		[Como Contexto]
		[Como Habilidad]
		[Como Hacer Matemáticas]
		]
	\end{forest}
	
	\subsection{Desglose de ``Como Contexto''}
	
	\begin{enumerate}
		\item Justificación para enseñar matemáticas.
		\item Motivación para ciertos temas.
		\item Actividad recreativa.
		\item Desarrollo de habilidades.
		\item Práctica.
	\end{enumerate}
	
	\section{Taxonomía 14: Tipos de Procedimientos para Resolver Problemas}
	
	\begin{forest}
		for tree={draw,rounded corners}
		[Procedimientos
		[Algorítmicos]
		[Heurísticos
		[Variación]
		[Generalización]
		[Particularización]
		[Analogía]
		]
		]
	\end{forest}
	
	\section{Taxonomía 15: Objetivos de la Resolución de Problemas}
	
	\begin{forest}
		for tree={draw,rounded corners}
		[Objetivos
		[Objetivo Metodológico]
		[Objetivo Cognitivo]
		]
	\end{forest}
	
	\section{Taxonomía 16: Conocimientos Necesarios para Resolver Problemas}
	
	\subsection{Nota}
	
	En el fragmento analizado aparece explícitamente:
	
	\begin{itemize}
		\item Conocimiento Lingüístico y Semántico.
	\end{itemize}
	
	Se identifican referencias a una clasificación más amplia que continúa en secciones posteriores del documento. La estructura completa no puede reconstruirse con certeza a partir del contenido analizado, por lo que se marca como clasificación parcialmente visible.
	
	\section{Taxonomía 17: Campos que Nutren la Educación Matemática}
	
	\begin{forest}
		for tree={draw,rounded corners}
		[Educación Matemática
		[Pedagogía]
		[Didáctica]
		[Psicología]
		[Lingüística (Semiótica)]
		[Epistemología]
		[Sociología]
		[Ciencias Cognitivas]
		[Informática]
		[Matemática]
		]
	\end{forest}
	
	\section{Mapa General de Relaciones}
	
	\begin{center}
		\begin{tikzpicture}[
			node distance=2.5cm,
			every node/.style={draw,rounded corners,align=center}
			]
			\node (sistemas) {Sistemas\Matemáticos};
			\node (pensamiento) [below left of=sistemas] {Pensamiento\Matemático};
			\node (subprocesos) [below right of=sistemas] {Subprocesos};
			\node (semiotica) [below of=sistemas] {Representaciones\Semióticas};
			\node (problemas) [below of=subprocesos] {Resolución\de Problemas};
			\node (razonamiento) [below of=pensamiento] {Razonamiento};
			\node (comunicacion) [below of=semiotica] {Comunicación};
			
			\draw[->] (sistemas)--(pensamiento);
			\draw[->] (sistemas)--(subprocesos);
			\draw[->] (semiotica)--(comunicacion);
			\draw[->] (semiotica)--(razonamiento);
			\draw[->] (subprocesos)--(problemas);
			\draw[->] (subprocesos)--(razonamiento);
			\draw[->] (subprocesos)--(comunicacion);
		\end{tikzpicture}
	\end{center}
	
	\section{Resumen Consolidado}
	
	\begin{itemize}
		\item El eje central del documento es el desarrollo del pensamiento matemático.
		\item Los sistemas matemáticos constituyen la estructura conceptual básica.
		\item Los subprocesos representan la dimensión curricular transversal.
		\item Las representaciones semióticas constituyen el mecanismo cognitivo fundamental de acceso a los objetos matemáticos.
		\item La modelación, comunicación, razonamiento y resolución de problemas funcionan como procesos integrados.
		\item Las progresiones por grados aparecen organizadas mediante etapas evolutivas.
		\item La teoría semiótica de Duval actúa como marco articulador de múltiples clasificaciones presentes en el documento.
	\end{itemize}
	
	\section{Conclusiones}
	
	El documento presenta una organización altamente estructurada alrededor del desarrollo del pensamiento matemático. Las taxonomías identificadas no operan de manera independiente sino como un sistema integrado donde los sistemas matemáticos, los subprocesos curriculares, las representaciones semióticas y las progresiones evolutivas se encuentran estrechamente articulados. La modelación matemática, la comunicación matemática, el razonamiento matemático y la resolución de problemas constituyen los principales organizadores curriculares y se apoyan transversalmente en los sistemas de representación semiótica y en los diferentes sistemas matemáticos.
	
\end{document}
